\documentclass[a4paper,12pt]{article}
\usepackage[T2A]{fontenc}
\usepackage[utf8]{inputenc}
\usepackage[russian]{babel}
\usepackage{amsmath,amssymb,mathtools}
\renewcommand{\familydefault}{\sfdefault}
\usepackage{geometry}
\geometry{left=24mm,right=24mm,top=26mm,bottom=25mm,headheight=24pt}
\usepackage{microtype}
\usepackage{tikz}
\usetikzlibrary{arrows.meta,positioning,shapes.geometric,calc}
\usepackage{booktabs,tabularx}
\usepackage{enumitem}
\usepackage{fancyhdr}
\usepackage{setspace}
\usepackage{xcolor}
\usepackage{hyperref}
\definecolor{accent}{HTML}{007B80}
\definecolor{darkaccent}{HTML}{00585B}
\definecolor{textgray}{HTML}{3A3D40}
\definecolor{muted}{HTML}{6F7478}
\hypersetup{colorlinks=true,linkcolor=darkaccent,urlcolor=darkaccent}
\onehalfspacing
\setlength{\parindent}{0pt}
\setlength{\parskip}{0.35em}
\setlist[itemize]{leftmargin=1.9em,itemsep=0.10em,topsep=0.25em}
\setlist[enumerate]{leftmargin=1.9em,itemsep=0.45em,topsep=0.25em}
\pagestyle{fancy}
\fancyhf{}
\fancyhead[L]{\footnotesize\color{muted} Движение по окружности}
\fancyhead[R]{\footnotesize\color{muted} Володя Хачатурян}
\fancyfoot[C]{\footnotesize\thepage}
\renewcommand{\headrulewidth}{0.45pt}
\renewcommand{\headrule}{\hbox to\headwidth{\color{accent}\leaders\hrule height \headrulewidth\hfill}}
\newcommand{\sechead}[1]{\vspace{0.35em}{\large\bfseries\color{darkaccent}#1}\par\vspace{0.22em}{\color{accent}\rule{\linewidth}{0.48pt}}\vspace{0.35em}}
\newcommand{\subhead}[1]{\vspace{0.32em}{\normalsize\bfseries\color{darkaccent}#1}\par}
\newcommand{\tip}[1]{{\color{darkaccent}\bfseries Важно.} #1}
\newcommand{\ex}[1]{{\color{darkaccent}\bfseries Пример #1.}\enspace}
\newcommand{\ac}{a_{\text{ц}}}
\newcommand{\ft}{F_{\text{тр}}}
\newcommand{\fmax}{F_{\text{тр},\max}}
\begin{document}

% ТИТУЛЬНЫЙ ЛИСТ
\begin{titlepage}
\thispagestyle{empty}
\centering
\vspace*{14mm}
{\large\color{darkaccent} ФИЗИКА\par}
\vspace{15mm}
{\Huge\color{textgray} Чек-лист\par}
\vspace{10mm}
{\LARGE\color{accent} Движение по окружности\par}
\vspace{3mm}
{\Large\color{accent} Центростремительная сила, трение\par}
{\Large\color{accent} и натяжение нити\par}
\vspace{13mm}
{\large\color{textgray} Второй закон Ньютона в проекциях\par}
\vfill
{\large Володя Хачатурян\par}
\vspace{9mm}
{\large 21 ноября 2026 г.\par}
\vspace{10mm}
{\normalsize Лёвин Артём Александрович\par}
{\normalsize эксклюзивно для Володи Хачатуряна\par}
\vspace{10mm}
\begin{tikzpicture}[x=1cm,y=1cm]
\draw[accent,line width=0.8pt] (0,0) .. controls (1.6,.48) and (2.35,-.48) .. (4.0,0)
.. controls (5.45,.47) and (6.45,-.48) .. (8.0,0)
.. controls (9.15,.34) and (10.15,-.35) .. (11.35,0);
\fill[darkaccent] (4.0,0) circle(1.4pt);
\end{tikzpicture}
\vspace{3mm}
\end{titlepage}

\sechead{1. Что означает движение по окружности}

При движении с постоянной \emph{по модулю} скоростью направление вектора скорости непрерывно меняется. Поэтому у тела есть \textbf{центростремительное ускорение}, направленное к центру окружности.

\begin{center}
\begin{tikzpicture}[scale=0.92,>=Latex]
\coordinate (O) at (0,0); \coordinate (P) at (2.0,0.95);
\draw[muted,densely dashed] (O) circle (2.21);
\draw[muted] (O) -- (P) node[midway,below] {$R$};
\fill[muted] (O) circle (1.4pt) node[below left] {$O$};
\fill[accent] (P) circle (3.3pt) node[above right] {$m$};
\draw[->,thick,accent] (P) -- ++(-1.6,-.76) node[pos=.63,below] {$\vec a_{\text{ц}}$};
\draw[->,thick,red!65!black] (P) -- ++(-.85,1.78) node[pos=.62,right] {$\vec v$};
\end{tikzpicture}
\end{center}

\subhead{Три равноценных формулы}
\[
\boxed{\ac=\frac{v^2}{R}=\omega^2R=\frac{4\pi^2R}{T^2}}.
\]
Здесь $R$ --- радиус окружности (м), $v$ --- линейная скорость (м/с), $\omega$ --- угловая скорость (рад/с), $T$ --- период одного оборота (с). В разделе о нити буква $T$ обозначает натяжение; эти величины различайте по смыслу условия.
\[
\boxed{v=\omega R=\frac{2\pi R}{T}},\qquad
\boxed{\omega=\frac{2\pi}{T}}.
\]

\subhead{Главное следствие второго закона Ньютона}
\[
\boxed{\sum F_{r}=m\ac=m\frac{v^2}{R}=m\omega^2R}.
\]
Ось $r$ направляем \emph{к центру}. Величина $\sum F_r$ --- сумма проекций \textbf{реальных сил} на эту ось. Словами «центростремительная сила» обозначают их равнодействующую в радиальном направлении.

\tip{При равномерном движении по окружности ускорение направлено к центру, хотя мгновенная скорость направлена по касательной.}

\clearpage
\sechead{2. Горизонтальная поверхность: трение удерживает тело}

Брусок движется по окружности на горизонтальном диске, не проскальзывая относительно него. Силы, действующие на брусок:
\begin{itemize}
\item $m\vec g$ --- вертикально вниз;
\item $\vec N$ --- реакция опоры, вертикально вверх;
\item $\vec F_{\text{тр}}$ --- сила трения покоя, горизонтально к оси вращения.
\end{itemize}

\begin{center}
\begin{tikzpicture}[>=Latex,scale=0.93]
\draw[muted,line width=.7pt] (-3.4,-.65)--(3.35,-.65);
\draw[fill=accent!10,draw=darkaccent,rounded corners=3pt] (-.47,-.64) rectangle (.47,-.04);
\coordinate (B) at (0,-.03);
\draw[->,very thick,accent] (B) -- +(0,1.6) node[above] {$\vec N$};
\draw[->,very thick,red!65!black] (B) -- +(0,-1.53) node[below] {$m\vec g$};
\draw[->,very thick,darkaccent] (B) -- +(2,0) node[right] {$\vec F_{\text{тр}}$};
\node[below,muted] at (2,-.67) {к центру};
\end{tikzpicture}
\end{center}

В вертикальном направлении ускорение отсутствует:
\[
\sum F_y=N-mg=0 \quad\Longrightarrow\quad \boxed{N=mg}.
\]
В направлении центра окружности:
\[
\boxed{\ft=m\ac=m\frac{v^2}{R}=m\omega^2R}.
\]

\subhead{Когда тело начинает скользить?}
\emph{Трение покоя подстраивается под требуемую силу} до предельного значения:
\[
\ft\le \fmax=\mu N=\mu mg.
\]
Для движения без проскальзывания необходимо
\[
 m\frac{v^2}{R}\le\mu mg
 \quad\Longrightarrow\quad
 \boxed{v\le v_{\max}=\sqrt{\mu gR}},\qquad
 \boxed{\omega\le\sqrt{\frac{\mu g}{R}}}.
\]
Здесь $\mu$ --- коэффициент трения покоя; $g$ --- ускорение свободного падения.

\tip{Равенство $\ft=\mu N$ выполняется на границе скольжения. При меньших скоростях сила трения обычно меньше своего максимума.}

\clearpage
\sechead{3. Два типовых расчёта с доски}

\subhead{Пример 1. Брусок на вращающемся диске}
\ex{1}Брусок массой $0{,}5$ кг лежит на расстоянии $0{,}4$ м от оси диска. Диск вращается с угловой скоростью $5$ рад/с. Брусок покоится относительно диска. Найдите силу трения.

\textbf{Дано:} $m=0{,}5$ кг, $R=0{,}4$ м, $\omega=5$ рад/с.

\textbf{Решение.} Сила трения сообщает бруску центростремительное ускорение:
\[
\ac=\omega^2R=5^2\cdot0{,}4=10\ \text{м/с}^2.
\]
\[
\ft=m\ac=0{,}5\cdot10=\boxed{5\ \text{Н}}.
\]
\textbf{Проверка размерности:} кг$\cdot$м/с$^2$ = Н.

\subhead{Пример 2. Автомобиль на горизонтальном повороте}
\ex{2}Автомобиль массой $1000$ кг проходит горизонтальный поворот радиусом $50$ м со скоростью $10$ м/с. Найдите силу трения, необходимую для движения по окружности.

\textbf{Дано:} $m=1000$ кг, $R=50$ м, $v=10$ м/с.

\textbf{Решение.} Здесь удобнее воспользоваться линейной скоростью:
\[
\ft=m\frac{v^2}{R}=1000\cdot\frac{10^2}{50}
=\boxed{2000\ \text{Н}}.
\]
Если принять $g=10$ м/с$^2$, минимальный коэффициент трения, допускающий такой поворот без скольжения, равен
\[
\mu_{\min}=\frac{\ft}{mg}=\frac{2000}{1000\cdot10}=0{,}2.
\]

\subhead{Как распознать нужную формулу}
\begin{center}
\renewcommand{\arraystretch}{1.28}
\begin{tabularx}{\linewidth}{@{}p{.37\linewidth}X@{}}
\toprule
\textbf{В условии дано} & \textbf{Удобно использовать}\\
\midrule
Радиус и линейная скорость & $\displaystyle\ac=\frac{v^2}{R}$\\
Радиус и угловая скорость & $\displaystyle\ac=\omega^2R$\\
Радиус и период & $\displaystyle\ac=\frac{4\pi^2R}{T^2}$\\
Коэффициент трения и радиус & $\displaystyle v_{\max}=\sqrt{\mu gR}$\\
\bottomrule
\end{tabularx}
\end{center}

\clearpage
\sechead{4. Тело на наклонной нити: разложение силы натяжения}

Тело движется по горизонтальной окружности. Нить образует угол $\alpha$ \textbf{с вертикалью}. Натяжение $\vec T$ направлено вдоль нити, а сила тяжести --- вниз.

\begin{center}
\begin{tikzpicture}[scale=.98,>=Latex]
\coordinate (A) at (0,2.8); \coordinate (B) at (1.9,0); \coordinate (C) at (0,0);
\draw[muted,densely dashed] (A)--(C);
\draw[darkaccent,line width=1.0pt] (A)--(B);
\draw[muted,densely dashed] (-2,0)--(3.3,0);
\draw[muted,densely dashed] (C) ellipse (1.9 and .36);
\fill[muted] (A) circle (2pt);
\fill[darkaccent] (B) circle (3.5pt);
\node[above right] at (A) {точка подвеса};
\node[below right] at (B) {$m$};
\draw[->,very thick,accent] (B)--($(B)!0.48!(A)$) node[above right] {$\vec T$};
\draw[->,very thick,red!65!black] (B)--++(0,-1.2) node[right] {$m\vec g$};
\draw[->,thick,darkaccent] (B)--++(-1.35,0) node[above] {$T_x$};
\draw[->,thick,darkaccent] (B)--++(0,1.3) node[right] {$T_y$};
\draw (0,2.05) arc[start angle=-90,end angle=-55.8,radius=.75];
\node at (.38,1.92) {$\alpha$};
\draw[<->,muted] (0,-.58)--(1.9,-.58) node[midway,below] {$R$};
\end{tikzpicture}
\end{center}

\subhead{Разложение натяжения при угле от вертикали}
\[
\boxed{T_x=T\sin\alpha},\qquad \boxed{T_y=T\cos\alpha}.
\]
По вертикали тело остаётся на одной высоте; по горизонтали направлено центростремительное ускорение:
\[
\begin{cases}
T\sin\alpha=m\dfrac{v^2}{R}=m\omega^2R,\\[3pt]
T\cos\alpha=mg.
\end{cases}
\]
Отсюда
\[
\boxed{\tan\alpha=\frac{v^2}{gR}=\frac{\omega^2R}{g}},
\qquad \boxed{T=\frac{mg}{\cos\alpha}}.
\]

\tip{Если угол $\beta$ отсчитывают от горизонтали, проекции меняются местами: $T_x=T\cos\beta$, $T_y=T\sin\beta$. Перед вычислением всегда отмечайте, от какой оси задан угол.}

\clearpage
\sechead{5. Числовой пример для наклонной нити}
\ex{3}Тело массой $0{,}2$ кг движется по горизонтальной окружности радиусом $1$ м с угловой скоростью $\sqrt{10}$ рад/с. Нить составляет угол $\alpha$ с вертикалью. Примите $g=10$ м/с$^2$. Найдите натяжение и угол.

\textbf{Шаг 1.} Найдём горизонтальную проекцию силы натяжения:
\[
T\sin\alpha=m\omega^2R=0{,}2\cdot10\cdot1=2\ \text{Н}.
\]
\textbf{Шаг 2.} Вертикальная проекция уравновешивает тяжесть:
\[
T\cos\alpha=mg=0{,}2\cdot10=2\ \text{Н}.
\]
\textbf{Шаг 3.} Делим первое равенство на второе:
\[
\tan\alpha=\frac{2}{2}=1\quad\Longrightarrow\quad \boxed{\alpha=45^\circ}.
\]
\textbf{Шаг 4.} Вычислим натяжение по вертикальной проекции:
\[
T=\frac{2}{\cos45^\circ}=\boxed{2\sqrt2\ \text{Н}}.
\]
\textbf{Проверка:} $T\sin45^\circ=2\sqrt2\cdot\sqrt2/2=2$ Н.

\subhead{Геометрическая связь с радиусом}
Если длина нити $\ell$, а её верхняя точка лежит на оси вращения, то
\[
\boxed{R=\ell\sin\alpha}\qquad
\text{(угол от вертикали).}
\]
Эта формула получается из прямоугольного треугольника: радиус --- катет, противоположный углу $\alpha$, а нить --- гипотенуза.

\subhead{Три проверки перед окончательным ответом}
\begin{itemize}
\item Направлена ли радиальная проекция силы к центру окружности?
\item Верно ли определён угол: к вертикали или к горизонтали?
\item Можно ли проверить результат из второго уравнения проекций?
\end{itemize}

\clearpage
\begin{center}
{\Large\bfseries\color{darkaccent} Алгоритм: второй закон Ньютона при движении по окружности}\par
\vspace{1.2cm}
\begin{tikzpicture}[
 arr/.style={-{Latex[length=2.6mm]},thick,draw=darkaccent},
 term/.style={ellipse,draw=darkaccent,thick,minimum height=1.1cm,minimum width=4.0cm,align=center},
 act/.style={rectangle,rounded corners=3pt,draw=accent,thick,minimum height=1.25cm,text width=9cm,align=center,inner sep=6pt},
 choose/.style={diamond,aspect=2.1,draw=darkaccent,thick,align=center,text width=3.7cm,inner sep=2pt},
 branch/.style={rectangle,rounded corners=3pt,draw=accent,thick,text width=5.2cm,minimum height=2.0cm,align=center,inner sep=5pt},
 font=\small
]
\node[term] (start) at (0,0) {Начало задачи};
\node[act] (axis) at (0,-2) {Найдите центр окружности. Радиальную ось направьте к центру, вертикальную --- вверх.};
\node[act] (forces) at (0,-4.3) {Изобразите реальные силы. Запишите $\sum F_y=0$ и $\sum F_r=m\ac$.};
\node[choose] (kind) at (0,-7.15) {Какая сила даёт радиальную проекцию?};
\node[branch] (fric) at (-4.2,-10.35) {\textbf{Трение}\par $\ft=m\ac$;\par $N=mg$; при проверке сцепления $\ft\le\mu N$};
\node[branch] (rope) at (4.2,-10.35) {\textbf{Натяжение}\par $T\sin\alpha=m\ac$;\par $T\cos\alpha=mg$ (угол от вертикали)};
\node[act] (calc) at (0,-13.2) {Подставьте $\ac=v^2/R=\omega^2R$. Найдите искомое, проверьте единицы и ограничения.};
\node[term] (end) at (0,-15.4) {Ответ};
\draw[arr] (start)--(axis);
\draw[arr] (axis)--(forces);
\draw[arr] (forces)--(kind);
\draw[arr] (kind.west)-| node[pos=.29,above,sloped]{трение} (fric.north);
\draw[arr] (kind.east)-| node[pos=.29,above,sloped]{нить} (rope.north);
\draw[arr] (fric.south)|-(calc.west);
\draw[arr] (rope.south)|-(calc.east);
\draw[arr] (calc)--(end);
\end{tikzpicture}
\end{center}
\clearpage

\sechead{6. Частые ошибки и короткая памятка}
\begin{enumerate}[label=\arabic*.]
\item \textbf{Направляют центростремительную силу по касательной.} По касательной направлен вектор скорости; радиальное ускорение всегда направлено к центру.
\item \textbf{Добавляют «центростремительную силу» поверх уже изображённых сил.} Радиальная равнодействующая складывается из реально действующих сил: трения, натяжения и т.~д.
\item \textbf{Приравнивают трение к $\mu N$ при любой скорости.} Для трения покоя это верхняя граница, достижимая на пороге скольжения.
\item \textbf{Путают формулы $v^2/R$ и $\omega^2R$.} Проверяйте, какая скорость дана: в м/с или рад/с.
\item \textbf{Меняют синус и косинус местами.} Сначала отметьте, от какой оси измеряется угол наклона нити.
\item \textbf{Смешивают радиус, длину нити и диаметр.} Во всех формулах ускорения $R$ --- расстояние до оси вращения.
\end{enumerate}

\subhead{Проверьте себя}
\begin{itemize}[label=$\square$]
\item Я понимаю, почему ускорение направлено к центру при постоянном модуле скорости.
\item Я умею переходить между $v$, $\omega$ и $T$.
\item Я выделяю реальные силы и записываю второй закон Ньютона в проекциях.
\item Я различаю необходимую силу трения и максимальную силу трения покоя.
\item Я раскладываю натяжение нити по осям и учитываю, откуда отсчитывается угол.
\end{itemize}

\subhead{Сводная таблица}
\begin{center}
\renewcommand{\arraystretch}{1.3}
\begin{tabularx}{\linewidth}{@{}p{.36\linewidth}X@{}}
\toprule
\textbf{Ситуация} & \textbf{Ключевое уравнение}\\
\midrule
Любая окружность & $\sum F_r=mv^2/R=m\omega^2R$\\
Диск / горизонтальный поворот & $\ft=mv^2/R$, $N=mg$\\
Предел трения & $mv^2/R\le\mu mg$\\
Нить, $\alpha$ от вертикали & $T\sin\alpha=mv^2/R$, $T\cos\alpha=mg$\\
\bottomrule
\end{tabularx}
\end{center}

\clearpage
\sechead{Тренировочный блок --- 10 заданий}
Выполняйте задания последовательно. Если специально не оговорено иное, принимайте $g=10$ м/с$^2$. Промежуточные вычисления записывайте с единицами измерения.

\subhead{Базовый уровень}
\begin{enumerate}[label=\textbf{\arabic*.},leftmargin=2.1em]
\item Точка вращается по окружности радиусом $0{,}5$ м с угловой скоростью $4$ рад/с. Найдите линейную скорость и центростремительное ускорение.
\item Тело массой $0{,}25$ кг движется по окружности радиусом $2$ м со скоростью $6$ м/с. Найдите центростремительное ускорение и модуль радиальной равнодействующей сил.
\item Точка совершает полный оборот за $2$ с по окружности радиусом $1$ м. Найдите линейную скорость и центростремительное ускорение; ответ выразите через $\pi$.
\end{enumerate}

\subhead{Средний уровень}
\begin{enumerate}[label=\textbf{\arabic*.},start=4,leftmargin=2.1em]
\item Автомобиль массой $900$ кг проходит горизонтальный поворот радиусом $45$ м со скоростью $15$ м/с. Найдите необходимую силу трения между колёсами и дорогой.
\item Брусок массой $0{,}4$ кг удерживается на вращающемся горизонтальном диске на расстоянии $0{,}5$ м от оси. Угловая скорость диска $6$ рад/с. Найдите силу трения, действующую на брусок.
\end{enumerate}
\clearpage
\sechead{Тренировочный блок --- продолжение}
\begin{enumerate}[label=\textbf{\arabic*.},start=6,leftmargin=2.1em]
\item Коэффициент трения покоя между шинами и горизонтальной дорогой равен $0{,}45$. Найдите наибольшую скорость безопасного прохождения поворота радиусом $2$ м в данной идеализированной модели.
\item Брусок массой $0{,}6$ кг лежит на горизонтальном вращающемся диске на расстоянии $0{,}5$ м от оси. Угловая скорость $4$ рад/с, коэффициент трения покоя $0{,}5$. Может ли брусок вращаться вместе с диском без проскальзывания? Обоснуйте расчётом.
\end{enumerate}

\subhead{Повышенный уровень}
\begin{enumerate}[label=\textbf{\arabic*.},start=8,leftmargin=2.1em]
\item Тело массой $0{,}6$ кг движется по горизонтальной окружности на нити, отклонённой от \emph{вертикали} на $60^\circ$. Найдите силу натяжения нити и её горизонтальную составляющую.
\item Тело массой $0{,}5$ кг движется на наклонной нити по горизонтальной окружности радиусом $0{,}8$ м со скоростью $4$ м/с. Найдите натяжение нити и $\tan\alpha$, где $\alpha$ --- угол нити с вертикалью.
\item Брусок вращается вместе с горизонтальным диском на расстоянии $0{,}8$ м от оси. Коэффициент трения покоя $0{,}5$. Найдите наименьший период обращения, при котором проскальзывания ещё нет.
\end{enumerate}

\clearpage
\sechead{Ответы к тренировочному блоку}
\begin{center}
\renewcommand{\arraystretch}{1.42}
\begin{tabularx}{\linewidth}{@{}p{1.2cm}X@{}}
\toprule
\textbf{№} & \textbf{Ответ}\\
\midrule
1 & $v=2$ м/с; $\ac=8$ м/с$^2$.\\
2 & $\ac=18$ м/с$^2$; $F_r=4{,}5$ Н.\\
3 & $v=\pi$ м/с; $\ac=\pi^2$ м/с$^2$.\\
4 & $\ft=4500$ Н.\\
5 & $\ft=7{,}2$ Н.\\
6 & $v_{\max}=3$ м/с.\\
7 & Нет: требуется $4{,}8$ Н, доступно не более $3$ Н.\\
8 & $T=12$ Н; $T_x=6\sqrt3$ Н.\\
9 & $T=5\sqrt5$ Н; $\tan\alpha=2$.\\
10 & $T_{\min}=0{,}8\pi$ с.\\
\bottomrule
\end{tabularx}
\end{center}

\vspace{1.0em}
\subhead{Мини-памятка перед проверкой решения}
\begin{enumerate}[label=\arabic*.]
\item Постройте схему сил и укажите направление к центру.
\item Подберите формулу для ускорения по известным величинам.
\item Запишите второй закон Ньютона в радиальной и вертикальной проекциях.
\item Проверьте ограничения, единицы и физический смысл результата.
\end{enumerate}

\vfill
\begin{center}
{\footnotesize\color{muted} Лёвин Артём Александрович --- эксклюзивно для Володи Хачатуряна}
\end{center}

\end{document}
